\documentclass{beamer}[10]
\usepackage{pgf}
\usepackage{beamerthemesplit}

% Basado en: http://www.tjansson.dk/?p=82

\usepackage[utf8]{inputenc}


\definecolor{kugreen}{RGB}{149,122,143} %{50,93,61}
\definecolor{kugreenlys}{RGB}{132,158,139}
\definecolor{kugreenlyslys}{RGB}{173,190,177}
\definecolor{kugreenlyslyslys}{RGB}{214,223,216}
\setbeamercovered{transparent}
\mode<presentation>
\usetheme[numbers,totalnumber,compress,sidebarshades]{PaloAlto}

% Con número de página, o con footbar
%\setbeamertemplate{footline}[frame number]
\setbeamertemplate{footline}[frame]

  \usecolortheme[named=kugreen]{structure}
  \useinnertheme{circles}
  \usefonttheme[onlymath]{serif}
  \setbeamercovered{transparent}
  \setbeamertemplate{blocks}[rounded][shadow=true]
  \logo{\includegraphics[width=1.5cm]{imagenes/linti.pdf}}

% Ocultar barra de navegación
\beamertemplatenavigationsymbolsempty

\title{Programando con Python y Robots}
%\subtitle{la seguridad, la adecuada instalación y la libertad de elección}
\author{Fernando López}
\institute{LINTI \\ Facultad de Informática \\ Universidad Nacional de la Plata
}
\date{29 de Agosto de 2013}


\begin{document}

\frame{\titlepage \vspace{-0.5cm}

}


\frame
{
\frametitle{Indice}
\tableofcontents[pausesection]
}

\begin{frame}{Los inicios del proyectos}
\end{frame}
\begin{frame}{Pasantías}
\end{frame}
\begin{frame}{El proyecto en las escuelas}
\end{frame}
\begin{frame}{Software necesario}
	\begin{itemize}
		\item Duinobot $\rightarrow$ Lihuen/GitHub/Pip
		\item Firmata $\rightarrow$ GitHub/Preinstalado
	\end{itemize}
	\begin{itemize}
		\item Duinobot es la API en Python con la que estaremos en contacto.
		\item Duinobot incluye una copia modificada de PyFirmata encargada de la comunicación.
		\item Interfaz serial.
		\item El robot ejecuta una copia modificada de Firmata.
	\end{itemize}
\end{frame}
\begin{frame}{Primeros pasos}
	\begin{itemize}
		\item Libro.
		\item Intérprete.
		\item Conexión
		\item Beep.
		\item Movimientos.
		\item Movimientos simples.
		\item Figuras.
		\item Sensores y condicionales.
		\item Control remoto con raw\_input.
	\end{itemize}
\end{frame}
\begin{frame}{Algoritmos simples}
	\begin{itemize}
		\item Detenerse en obstáculo.
		\item Esquivar.
		\item Escapar.
		\item Mantenerse en el perímetro.
		\item Seguir línea.
		\item Beep proporcional a la distancia.
		\item Beep inversamente proporcional.
	\end{itemize}
\end{frame}
\frame
{
\frametitle{Direcciones}

\begin{block}{Direcciones de contacto}
\begin{center}
flopez@linti.unlp.edu.ar \\
soportelihuen@linti.unlp.edu.ar
\end{center}
\end{block}

}


%~ \frame
%~ {
%~ \frametitle{Entorno de pre-boot del BIOS (3)}
	%~ Hoy en día sería deseable:
	%~ \begin{itemize}
		%~ \item Acceso al stack TCP, por ejemplo para gestores de 
		%~ arranque por red (en lugar de depender de la option ROM con 
		%~ soporte PXE que pudiera tener la NIC).
		%~ \item Uso de dispositivos como un mouse.
		%~ \item La posibilidad de ampliar o mejorar el soporte de 
		%~ hardware sin option ROMs.
		%~ \item Acceso a direccionamiento de 32 o 64 bits.
	%~ \end{itemize}
%~ }
%~ 
%~ \frame{
%~ 
%~ Efteråret 2004
%~ \subsection{Valg af projekt og vejleder}
%~ \frametitle{Tilfældighed}
%~ 
%~ \begin{block}{Kryptering med kaotiske kredsløb}
%~ \begin{columns}
%~ \column{.3\textwidth} \hspace{0.5cm}
%~ \includegraphics[width=0.7\textwidth]{imagenes/linti-f} 
%~ \column{.7\textwidth}
%~ \textit{Mogens Høgh Jensen}, NBI
%~ \end{columns}
%~ \end{block}
%~ 
%~ \begin{block}{Sandriller}
%~ \begin{columns}
%~ \column{.3\textwidth} \hspace{0.5cm}
%~ \includegraphics[width=0.7\textwidth]{imagenes/linti-f} 
%~ \column{.7\textwidth}
%~ \textit{Clive Ellegaard}, NBI og \emph{Tomas Bohr}, DTU
%~ \end{columns}
%~ \end{block}
%~ }





\end{document}
